\documentclass[12pt,a4paper]{article}
\usepackage[margin=25mm, headheight=15pt, headsep=10pt]{geometry}
\usepackage{fontspec}
\usepackage[table]{xcolor} % Note: table option is required for rowcolors
\usepackage{titlesec}
\usepackage{tocloft}
\usepackage{fancyhdr}
\usepackage{booktabs}
\usepackage{tcolorbox}
\tcbuselibrary{skins, breakable}
\usepackage[colorlinks=true, linkcolor=emerald, urlcolor=emerald, bookmarks=true]{hyperref}

\definecolor{emerald}{RGB}{0,102,51}
\definecolor{darkred}{RGB}{139,0,0}
\definecolor{lightgray}{RGB}{245,245,245}

\usepackage[nil]{babel}
\babelprovide[import, main]{bengali}
\babelprovide[import]{arabic}
\babelfont{rm}[Renderer=HarfBuzz,Scale=1.0]{Noto Serif Bengali}
\babelfont{sf}[Renderer=HarfBuzz]{Noto Sans Bengali}
\babelfont{tt}[Renderer=HarfBuzz]{Noto Sans Bengali}
\babelfont[arabic]{rm}[Renderer=HarfBuzz,Scale=1.1]{Noto Naskh Arabic}

% Section styling
\titleformat{\section}{\Large\bfseries\color{emerald}}{\thesection.}{0.5em}{}
\titleformat{\subsection}{\large\bfseries\color{emerald!85!black}}{\thesubsection.}{0.5em}{}
\titleformat{\subsubsection}{\normalsize\bfseries\color{emerald!70!black}}{\thesubsubsection.}{0.5em}{}
\titlespacing*{\section}{0pt}{18pt}{8pt}

% Fancy Header/Footer setup
\pagestyle{fancy}
\fancyhf{} % clear all header and footer fields
\fancyhead[L]{\color{emerald}\small\textbf{\nouppercase{\leftmark}}}
\fancyfoot[L]{\scriptsize\color{black!50}{\lat{AI}}-সংকলিত, আলেম-যাচাইকৃত নয়}
\fancyfoot[C]{\color{emerald!70!black}\thepage}
\renewcommand{\headrulewidth}{0.4pt}
\renewcommand{\headrule}{\hbox to\headwidth{\color{emerald}\leaders\hrule height \headrulewidth\hfill}}
\renewcommand{\footrulewidth}{0pt}

% Custom boxes
% 1. Ayat/Hadith Box
\newtcolorbox{ayatbox}[1][]{
    enhanced, breakable, 
    colback=emerald!5, 
    colframe=emerald, 
    boxrule=0.5pt, 
    arc=3pt, 
    left=10pt, right=10pt, top=10pt, bottom=10pt,
    #1
}

% 2. Quote/Translation Box
\newtcolorbox{quotebox}[1][]{
    enhanced, breakable,
    frame hidden,
    colback=lightgray,
    borderline west={3pt}{0pt}{emerald},
    left=15pt, right=10pt, top=10pt, bottom=10pt,
    sharp corners,
    #1
}

% 3. AI-generated notice box
\newtcolorbox{warnbox}[1][]{
    enhanced, breakable,
    colback=red!4,
    colframe=darkred,
    boxrule=0.8pt,
    arc=3pt,
    left=10pt, right=10pt, top=10pt, bottom=10pt,
    #1
}

% Commands
\newcommand{\arab}[1]{\foreignlanguage{arabic}{#1}}
\newcommand{\kib}[1]{\textcolor{emerald}{\textbf{#1}}}

% Latin (English) fallback: Noto Serif Bengali has NO Latin glyphs
\newfontfamily\latfont[Renderer=HarfBuzz]{Noto Serif}
\newcommand{\lat}[1]{{\latfont #1}}

\setlength{\parskip}{5pt}
\setlength{\parindent}{0pt}

% Fix for missing bullet in Noto Serif Bengali
\renewcommand{\labelitemi}{$\bullet$}



\begin{document}
\begin{center}{\Large\bfseries\color{emerald} কুরআনের আয়াত — সালাত (নামাজ)}\end{center}
\vspace{1em}
\section*{কুরআনের আয়াত — সালাত (নামাজ)}

\begin{warnbox}
 \textbf{গুরুত্বপূর্ণ নোটিশ (\lat{Important} \lat{Notice}):} এই কন্টেন্টটি \lat{AI} (কৃত্রিম বুদ্ধিমত্তা) দিয়ে প্রস্তুত ও সংকলিত। \lat{AI} কঠোর সূত্র-মান নীতি মেনে শুধু নির্ভরযোগ্য উৎস থেকে তথ্য সংগ্রহ করেছে — কেবল সহীহ/হাসান হাদিস ও সহীহ সনদের আসার রাখা হয়েছে, দুর্বল সনদ বাদ দেওয়া হয়েছে। তবে এটি কোনো আলেম বা মুহাদ্দিস কর্তৃক যাচাইকৃত নয়।
\end{warnbox}


\begin{quote}
\textbf{পড়ার নিয়ম:} প্রতিটি আয়াতের সাথে থাকবে — \textbf{মূল আরবি} $\rightarrow$ \textbf{বাংলা উচ্চারণ} $\rightarrow$ \textbf{বাংলা অর্থ (\lat{pure})} $\rightarrow$ \textbf{শব্দের ব্যাখ্যা} $\rightarrow$ \textbf{গভীর তাফসীর} $\rightarrow$ \textbf{শিক্ষা ও প্রয়োগ}। \\
তাফসীর সংগ্রহ করা হয়েছে প্রামাণ্য গ্রন্থ থেকে: তাফসীর ইবনে কাসীর, জামিউল বায়ান (তাবারী), তাফসীরুল কুরতুবী, তাফসীর আস-সা'দী, মাআরিফুল কুরআন (মুফতী মুহাম্মদ শাফী উসমানী)। আয়াতের আরবি ও অনুবাদ যাচাই করা হয়েছে কুরআনের প্রামাণ্য পাঠ ও অনুবাদ থেকে।
\end{quote}


\medskip\hrule\medskip

\subsection*{আয়াত ১ — সূরা ত্বহা (২০) — আয়াত ১৪}

\subsubsection*{১. সূত্র কার্ড}

\begin{table}[h]\centering\small
\begin{tabular}{ll}
বিষয় & বিবরণ \\
\hline
\textbf{সূরা} & ত্বহা (২০) — মক্কী \\
\textbf{আয়াত} & ১৪ \\
\textbf{পারা} & ১৬ \\
\textbf{মূল বিষয়} & নামাজ কায়েমের মূল উদ্দেশ্য — আল্লাহর স্মরণ (যিকর) \\
\textbf{বক্তা} & আল্লাহ — মুসা (আঃ)-কে সম্বোধন করে \\
\hline
\end{tabular}
\end{table}

\subsubsection*{২. প্রসঙ্গ ও প্রেক্ষাপট (\lat{Context})}

পবিত্র উপত্যকা তুয়ায় আল্লাহ মুসা (আঃ)-এর সাথে কথা বলেছেন। আয়াত ১৪-এ সেই সংলাপের কেন্দ্রীয় নির্দেশ — \lat{“}নিশ্চয়ই আমিই আল্লাহ। আমার ইবাদত করো এবং আমার স্মরণের জন্য নামাজ কায়েম করো।\lat{”} এখানে নামাজকে \textbf{স্মরণের মাধ্যম (\lat{means} \lat{of} \lat{remembrance})} বলা হয়েছে — অর্থাৎ নামাজের গভীরতম অর্থ হলো আল্লাহর সাথে সংযোগ, শুধু আনুষ্ঠানিকতা নয়।

\subsubsection*{৩. মূল আরবি}

\begin{center}\arab{إِنَّنِي أَنَا اللَّهُ لَا إِلَٰهَ إِلَّا أَنَا فَاعْبُدْنِي وَأَقِمِ الصَّلَاةَ لِذِكْرِي}\end{center}

\textbf{উচ্চারণ:} ইন্নানী আনা-ল্লা-হু লা- ইলা-হা ইল্লা- আনা- ফা'বুদনী ওয়া আক্বিমিস-সালা-তা লিযিকরী।

\subsubsection*{৪. বাংলা অর্থ (\lat{pure})}

\textbf{\lat{“}নিশ্চয়ই আমিই আল্লাহ, আমি ছাড়া কোনো উপাস্য নেই। অতএব আমারই ইবাদত করো এবং আমার স্মরণের জন্য নামাজ কায়েম করো।\lat{”}}

\subsubsection*{৫. শব্দের ব্যাখ্যা (\lat{Vocabulary})}

\begin{table}[h]\centering\small
\begin{tabular}{ll}
শব্দ & অর্থ \\
\hline
\textbf{\arab{فَاعْبُدْنِي}} & আমারই ইবাদত করো — ইবাদত একমাত্র আল্লাহর জন্য \\
\textbf{\arab{وَأَقِمِ}} & কায়েম করো — প্রতিষ্ঠা করো, স্থায়ীভাবে পালন করো \\
\textbf{\arab{الصَّلَاةَ}} & নামাজ — এখানে শব্দটির মূল অর্থ দোয়া/সংযোগ \\
\textbf{\arab{لِذِكْرِي}} & আমার স্মরণের জন্য — আমার যিকরের উদ্দেশ্যে \\
\hline
\end{tabular}
\end{table}

\subsubsection*{৬. গভীর তাফসীর (\lat{Deep} \lat{Tafsir})}

\begin{itemize}
\item \textbf{ইবনে কাসীর:} আয়াতের অর্থ — আমার স্মরণ ও আমার প্রতি ভালোবাসার জন্য নামাজ কায়েম করো। কারো কারো মতে \lat{“}লিযিকরী\lat{”} অর্থ — যে নামাজ ভুলে যায়, তাকে স্মরণ করিয়ে দিতে নামাজ কায়েম করো; অথবা নামাজ পড়ো যাতে আমি তোমাকে স্মরণ করি। সব অর্থই পরস্পর সম্পূরক। (তাফসীর ইবনে কাসীর, সূরা ত্বহা)
\item \textbf{তাবারী:} আল্লাহ মুসা (আঃ)-কে আদেশ দিচ্ছেন — আমার ইবাদতের জন্য নামাজ প্রতিষ্ঠা করো, আমার স্মরণের জন্যই নামাজ আদায় করো। নামাজের অন্তর্নিহিত লক্ষ্য আল্লাহর স্মরণ। (জামিউল বায়ান, সূরা ত্বহা)
\item \textbf{সা'দী:} নামাজ যেন কেবল রুকন-সিজদার \textbf{\lat{routine} (অভ্যাস)} না হয়; নামাজের প্রাণ হলো আল্লাহর স্মরণ — অন্তর আল্লাহর উপস্থিতি অনুভব করবে। (আস-সা'দী, সূরা ত্বহা)
\item \textbf{মাআরিফুল কুরআন:} আয়াত থেকে জানা যায় — নামাজের সঠিক আদায় আল্লাহর স্মরণকে গভীর করে; যে নামাজে আল্লাহর স্মরণ নেই, তা আত্মার খোরাক হয় না। (মাআরিফুল কুরআন, সূরা ত্বহা)
\end{itemize}
\subsubsection*{৭. সংশ্লিষ্ট হাদিস}

\begin{itemize}
\item \textbf{হাদিস (বুখারি ৭৫৬, মুসলিম ৩৯৪):} \lat{“}যে ব্যক্তি সূরা ফাতিহা পড়েনি, তার নামাজ নেই।\lat{”} — নামাজ কায়েম মানে ফাতিহাসহ পূর্ণ নামাজ।
\item \textbf{হাদিস (মুসলিম ৪৮২):} \lat{“}বান্দা তার রবের সবচেয়ে নিকটবর্তী হয় সিজদার সময়, তাই সিজদায় বেশি দোয়া করো।\lat{”} — নামাজই আল্লাহর সাথে নৈকট্যের মাধ্যম।
\end{itemize}
\subsubsection*{৮. গভীর শিক্ষা (\lat{Deep} \lat{Lessons})}

\begin{enumerate}
\item \textbf{নামাজের লক্ষ্য — স্মরণ।} নামাজ পড়া মানে আল্লাহকে ভুলে থাকা নয়; বরং সারাদিনের গাফিলতির মধ্যে নামাজই আল্লাহর স্মরণের \textbf{\lat{anchor} (নোঙর)}।
\item \textbf{ইবাদত আল্লাহর জন্যই।} \lat{“}\arab{فَاعْبُدْنِي}\lat{”} — ইবাদতের একমাত্র লক্ষ্য আল্লাহ; নামাজে রিয়া (প্রদর্শন) এই লক্ষ্যের বিপরীত।
\item \textbf{কায়েম শব্দের গভীরতা।} \lat{“}\arab{أَقِمِ}\lat{”} মানে শুধু পড়া নয় — প্রতিষ্ঠা করা: সময়মতো, জামাতে, নিয়মিত।
\end{enumerate}
\subsubsection*{৯. জীবনে প্রয়োগ (\lat{Practical} \lat{Application})}

\begin{enumerate}
\item নামাজে দাঁড়ানোর আগে নিয়ত করুন — এই নামাজ আল্লাহর স্মরণের জন্য।
\item নামাজের পর কিছুক্ষণ আল্লাহর যিকরে বসুন — যেন নামাজের \textbf{\lat{impact} (প্রভাব)} দিনে ছড়িয়ে পড়ে।
\end{enumerate}
\medskip\hrule\medskip

\subsection*{আয়াত ২ — সূরা আল-বাকারাহ (২) — আয়াত ৪৩}

\subsubsection*{১. সূত্র কার্ড}

\begin{table}[h]\centering\small
\begin{tabular}{ll}
বিষয় & বিবরণ \\
\hline
\textbf{সূরা} & আল-বাকারাহ (২) — মাদানী \\
\textbf{আয়াত} & ৪৩ \\
\textbf{পারা} & ১ \\
\textbf{মূল বিষয়} & নামাজ ও জাকাত একসাথে — সামাজিক দ্বীন; জামাতের নির্দেশ \\
\textbf{বক্তা} & আল্লাহ — বনী ইসরাঈলকে সম্বোধন \\
\hline
\end{tabular}
\end{table}

\subsubsection*{২. প্রসঙ্গ ও প্রেক্ষাপট (\lat{Context})}

সূরার শুরুতে বনী ইসরাঈলের কাছে অঙ্গীকারের কথা বলা হয়েছে — তারা আল্লাহর ইবাদত করবে, পিতামাতার সেবা করবে, নামাজ কায়েম করবে। আয়াত ৪৩-এ নির্দেশ এসেছে: নামাজ কায়েম করো, জাকাত দাও এবং \textbf{রুকুকারীদের সাথে রুকু করো} — অর্থাৎ জামাতের সাথে নামাজ পড়ো। এটি নামাজের সামাজিক দিক — একা নয়, উম্মতের সাথে।

\subsubsection*{৩. মূল আরবি}

\begin{center}\arab{وَأَقِيمُوا الصَّلَاةَ وَآتُوا الزَّكَاةَ وَارْكَعُوا مَعَ الرَّاكِعِينَ}\end{center}

\textbf{উচ্চারণ:} ওয়া আক্বীমুস-সালা-তা ওয়া আ-তুয-যাকা-তা ওয়ারকা'ঊ মা'আর-রা-কি'ঈন।

\subsubsection*{৪. বাংলা অর্থ (\lat{pure})}

\textbf{\lat{“}আর তোমরা নামাজ কায়েম করো, জাকাত প্রদান করো এবং রুকুকারীদের সাথে রুকু করো।\lat{”}}

\subsubsection*{৫. শব্দের ব্যাখ্যা (\lat{Vocabulary})}

\begin{table}[h]\centering\small
\begin{tabular}{ll}
শব্দ & অর্থ \\
\hline
\textbf{\arab{وَآتُوا}} & প্রদান করো — দাও \\
\textbf{\arab{الزَّكَاةَ}} & জাকাত — সম্পদের পবিত্রকরণ \\
\textbf{\arab{وَارْكَعُوا}} & রুকু করো — নামাজের রুকু \\
\textbf{\arab{مَعَ} \arab{الرَّاكِعِينَ}} & রুকুকারীদের সাথে — জামাতের সাথে \\
\hline
\end{tabular}
\end{table}

\subsubsection*{৬. গভীর তাফসীর (\lat{Deep} \lat{Tafsir})}

\begin{itemize}
\item \textbf{ইবনে কাসীর:} আয়াতে নামাজ-জাকাতকে একসাথে বলা হয়েছে — কারণ নামাজ আল্লাহর সাথে সম্পর্ক, জাকাত মানুষের সাথে সম্পর্ক; উভয়ের সমন্বয়ই পূর্ণ দ্বীন। \lat{“}রুকুকারীদের সাথে রুকু করো\lat{”} — জামাতে নামাজের নির্দেশ; আলেমগণ এটিকে জামাতের দলিল হিসেবে নিয়েছেন। (তাফসীর ইবনে কাসীর, সূরা বাকারাহ)
\item \textbf{কুরতুবী:} জামাতে নামাজ পড়ার আদেশ এখানে স্পষ্ট; সাহাবা ও তাবি'ঈনগণ এ আয়াত থেকে জামাতের গুরুত্ব বুঝেছেন। (তাফসীরুল কুরতুবী, সূরা বাকারাহ)
\item \textbf{মাআরিফুল কুরআন:} নামাজ ও জাকাতের \textbf{\lat{combination} (সম্মিলন)} ইসলামের মৌলিক কাঠামো; একটিকে বাদ দিয়ে অন্যটি পূর্ণ হয় না। (মাআরিফুল কুরআন, সূরা বাকারাহ)
\end{itemize}
\subsubsection*{৭. সংশ্লিষ্ট হাদিস}

\begin{itemize}
\item \textbf{হাদিস (বুখারি ৬৪৫, মুসলিম ৬৫০):} \lat{“}জামাতে নামাজ একাকী নামাজের চেয়ে ২৫/২৭ গুণ বেশি মর্যাদার।\lat{”}
\item \textbf{হাদিস (মুসলিম ৬৫৪):} \lat{“}যে ব্যক্তি ৪০ দিন জামাতে প্রথম তাকবীর পেল, তার জন্য দুইটি মুক্তি লেখা হয় — জাহান্নাম থেকে মুক্তি ও মুনাফিকি থেকে মুক্তি।\lat{”} (বর্ণনাটি মুসলিম ৬৫৪-এর সাথে সামঞ্জস্যপূর্ণ; আলেমরা এটিকে সহীহ বলেছেন।)
\end{itemize}
\subsubsection*{৮. গভীর শিক্ষা (\lat{Deep} \lat{Lessons})}

\begin{enumerate}
\item \textbf{নামাজ-জাকাত অবিচ্ছেদ্য।} আল্লাহর হক ও মানুষের হক — দুইটি একসাথে; এটাই দ্বীনের \textbf{\lat{balance} (সাম্য)}।
\item \textbf{জামাতের নির্দেশ।} আয়াতের শেষাংশ জামাতের দিকে ইঙ্গিত — নামাজ একা নয়, সমষ্টিগতভাবে কায়েম করতে হবে।
\item \textbf{আমলের স্থায়িত্ব।} \lat{“}কায়েম\lat{”} — শুধু একবার নয়, প্রতিষ্ঠা ও অব্যাহত রাখা।
\end{enumerate}
\subsubsection*{৯. জীবনে প্রয়োগ (\lat{Practical} \lat{Application})}

\begin{enumerate}
\item পাঁচ ওয়াক্ত জামাতে নামাজের \textbf{\lat{routine} (অভ্যাস)} গড়ুন — সম্ভব না হলে অন্তত ফজর ও ইশা।
\item নামাজের পাশাপাশি সম্পদের জাকাত/সদকা হিসাব করুন — দুইটি একসাথে রাখুন।
\end{enumerate}
\medskip\hrule\medskip

\subsection*{আয়াত ৩ — সূরা আন-নিসা (৪) — আয়াত ১০৩}

\subsubsection*{১. সূত্র কার্ড}

\begin{table}[h]\centering\small
\begin{tabular}{ll}
বিষয় & বিবরণ \\
\hline
\textbf{সূরা} & আন-নিসা (৪) — মাদানী \\
\textbf{আয়াত} & ১০৩ \\
\textbf{পারা} & ৫ \\
\textbf{মূল বিষয়} & নামাজ নির্দিষ্ট সময়ে ফরজ — সময়ের বাধ্যবাধকতা \\
\textbf{বক্তা} & আল্লাহ \\
\hline
\end{tabular}
\end{table}

\subsubsection*{২. প্রসঙ্গ ও প্রেক্ষাপট (\lat{Context})}

পূর্ববর্তী আয়াতে (৪:১০১-১০২) সফরে কসর ও যুদ্ধক্ষেত্রে (খাওফ) নামাজের নিয়ম এসেছে — এমনকি বিপদের মুহূর্তেও নামাজ ছাড়া যায় না। তারপর আয়াত ১০৩ ঘোষণা করে: নামাজ \textbf{নির্দিষ্ট সময়ে} মুমিনদের উপর ফরজ। অর্থাৎ নিরাপত্তা, সফর, যুদ্ধ — সব অবস্থায় নামাজ অপরিহার্য; শুধু পদ্ধতি পরিবর্তন হয়।

\subsubsection*{৩. মূল আরবি}

\begin{center}\arab{إِنَّ الصَّلَاةَ كَانَتْ عَلَى الْمُؤْمِنِينَ كِتَابًا مَوْقُوتًا}\end{center}

\textbf{উচ্চারণ:} ইন্নাস-সালা-তা কা-নাত 'আলাল মু'মিনীনা কিতা-বাম মাওকূতা-।

\subsubsection*{৪. বাংলা অর্থ (\lat{pure})}

\textbf{\lat{“}নিশ্চয়ই নামাজ নির্দিষ্ট সময়ে মুমিনদের উপর ফরজ।\lat{”}}

\subsubsection*{৫. শব্দের ব্যাখ্যা (\lat{Vocabulary})}

\begin{table}[h]\centering\small
\begin{tabular}{ll}
শব্দ & অর্থ \\
\hline
\textbf{\arab{كِتَابًا}} & লিখিত বিধান — ফরজ, নির্ধারিত \\
\textbf{\arab{مَوْقُوتًا}} & নির্দিষ্ট সময়ের — সময়-নির্ধারিত \\
\hline
\end{tabular}
\end{table}

\subsubsection*{৬. গভীর তাফসীর (\lat{Deep} \lat{Tafsir})}

\begin{itemize}
\item \textbf{ইবনে কাসীর:} আয়াতের অর্থ — নামাজের জন্য নির্দিষ্ট সময় নির্ধারিত, যেমন হজের জন্য নির্ধারিত সময়। সাহাবারা জিজ্ঞেস করতেন — কোন সময়ে কোন নামাজ; এ আয়াত সময়ের বাধ্যবাধকতা ঘোষণা করে। (তাফসীর ইবনে কাসীর, সূরা নিসা)
\item \textbf{সা'দী:} নামাজ এমন একটি ফরজ যা সময়ের সাথে \textbf{\lat{bound} (আবদ্ধ)} — সময়ের ভেতরে আদায় করা ওয়াজিব; সময় পার হওয়ার পর তা কাযা হয়। (আস-সা'দী, সূরা নিসা)
\item \textbf{মাআরিফুল কুরআন:} \lat{“}মাওকূত\lat{”} শব্দটি সময়-নির্ধারিত অর্থে; এতে নামাজের সময়মতো আদায়ের গুরুত্ব এবং সময়ের পূর্বে/পরে (বিনা কারণে) পড়া নিষিদ্ধ হওয়ার দিকে ইঙ্গিত। (মাআরিফুল কুরআন, সূরা নিসা)
\end{itemize}
\subsubsection*{৭. সংশ্লিষ্ট হাদিস}

\begin{itemize}
\item \textbf{হাদিস (মুসলিম ৬১০):} জিবরাইল (আঃ) নবীজিকে দুই দিন পাঁচ ওয়াক্ত নামাজের শুরু ও শেষ সময় শিখিয়েছেন — \lat{“}নামাজ এই দুই সময়ের মাঝে।\lat{”}
\item \textbf{হাদিস (বুখারি ৫৭৯, মুসলিম ৬০৮):} \lat{“}যে ব্যক্তি সূর্যাস্তের আগে আসরের এক রাকাত পেল, সে আসর পেল; যে সূর্যোদয়ের আগে ফজরের এক রাকাত পেল, সে ফজর পেল।\lat{”} — সময়ের শেষ প্রান্তে এক রাকাতও মূল্যবান।
\end{itemize}
\subsubsection*{৮. গভীর শিক্ষা (\lat{Deep} \lat{Lessons})}

\begin{enumerate}
\item \textbf{সময়ের আমানত।} নামাজের সময় আল্লাহর নির্ধারিত \textbf{\lat{boundary} (সীমা)} — সময়মতো নামাজ ঈমানের পরিচয়।
\item \textbf{সব অবস্থায় নামাজ।} যুদ্ধ-ভয়ের আয়াতের পর এই ঘোষণা — নামাজ কোনো অবস্থাতেই ছাড়ার নয়।
\item \textbf{কাযার গুরুত্ব।} সময় পার হলে নামাজ বাতিল হয় না, কাযা হয় — কিন্তু সময়মতো আদায়ই আসল।
\end{enumerate}
\subsubsection*{৯. জীবনে প্রয়োগ (\lat{Practical} \lat{Application})}

\begin{enumerate}
\item প্রতিটি নামাজ সময়মতো আদায়ের \textbf{\lat{schedule} (তালিকা)} মেনে চলুন — অজুহাত নয়।
\item নামাজের সময় হলে কাজ বন্ধ করে নামাজকে \textbf{\lat{priority} (অগ্রাধিকার)} দিন।
\end{enumerate}
\medskip\hrule\medskip

\subsection*{আয়াত ৪ — সূরা আল-বাকারাহ (২) — আয়াত ৪৫}

\subsubsection*{১. সূত্র কার্ড}

\begin{table}[h]\centering\small
\begin{tabular}{ll}
বিষয় & বিবরণ \\
\hline
\textbf{সূরা} & আল-বাকারাহ (২) — মাদানী \\
\textbf{আয়াত} & ৪৫ \\
\textbf{পারা} & ১ \\
\textbf{মূল বিষয়} & ধৈর্য ও নামাজের মাধ্যমে সাহায্য চাওয়া — ইস্তিআনাহ \\
\textbf{বক্তা} & আল্লাহ — বনী ইসরাঈলকে সম্বোধন \\
\hline
\end{tabular}
\end{table}

\subsubsection*{২. প্রসঙ্গ ও প্রেক্ষাপট (\lat{Context})}

বনী ইসরাঈলকে বলা হচ্ছে — বিপদে ধৈর্য (সবর) ও নামাজের সাহায্যে আল্লাহর কাছে সাহায্য চাও। এই আয়াতের সমান্তরাল বক্তব্য সূরা বাকারাহ ২:১৫৩-এ এসেছে মুমিনদের জন্য (\lat{“}তোমরা ধৈর্য ও নামাজের মাধ্যমে সাহায্য চাও\lat{”}) — যেখানে নামাজকে বিপদে \textbf{\lat{spiritual} \lat{support} (আধ্যাত্মিক সহায়তা)} বলা হয়েছে। নামাজ মানুষকে ধৈর্য, স্থিরতা ও আল্লাহর উপর ভরসা দেয়।

\subsubsection*{৩. মূল আরবি}

\begin{center}\arab{وَاسْتَعِينُوا بِالصَّبْرِ وَالصَّلَاةِ ۚ وَإِنَّهَا لَكَبِيرَةٌ إِلَّا عَلَى الْخَاشِعِينَ}\end{center}

\textbf{উচ্চারণ:} ওয়াসতা'ঈনূ বিস-সাবরি ওয়াস-সালা-হ, ওয়া ইন্নাহা- লাকাবীরাতুন ইল্লা- 'আলাল খা-শি'ঈন।

\subsubsection*{৪. বাংলা অর্থ (\lat{pure})}

\textbf{\lat{“}আর তোমরা ধৈর্য ও নামাজের মাধ্যমে সাহায্য চাও। নিশ্চয়ই এটি (নামাজ) কঠিন — কিন্তু বিনয়ী (খুশুকারী)দের জন্য সহজ।\lat{”}}

\subsubsection*{৫. শব্দের ব্যাখ্যা (\lat{Vocabulary})}

\begin{table}[h]\centering\small
\begin{tabular}{ll}
শব্দ & অর্থ \\
\hline
\textbf{\arab{وَاسْتَعِينُوا}} & সাহায্য চাও — ইস্তিআনাহ \\
\textbf{\arab{الصَّبْرِ}} & ধৈর্য — সবর \\
\textbf{\arab{الْخَاشِعِينَ}} & বিনয়ীরা — যারা নামাজে খুশু রাখে \\
\hline
\end{tabular}
\end{table}

\subsubsection*{৬. গভীর তাফসীর (\lat{Deep} \lat{Tafsir})}

\begin{itemize}
\item \textbf{ইবনে কাসীর:} নামাজকে এখানে সাহায্যের মাধ্যম বলা হয়েছে — কারণ নামাজ অন্তরকে আল্লাহর দিকে ফেরায়, বিপদে স্থির রাখে। \lat{“}কঠিন\lat{”} বলা হয়েছে — কারণ গাফিল অন্তরের জন্য নামাজ ভার; খুশুকারীর জন্য তা সহজ। (তাফসীর ইবনে কাসীর, সূরা বাকারাহ)
\item \textbf{তাবারী:} আয়াতের অর্থ — বিপদে ধৈর্য ধারণ করো এবং নামাজে আল্লাহর কাছে সাহায্য প্রার্থনা করো; নামাজই ধৈর্যের চর্চার শ্রেষ্ঠ মাধ্যম। (জামিউল বায়ান, সূরা বাকারাহ)
\item \textbf{সা'দী:} আয়াতে খুশুর গুরুত্ব স্পষ্ট — নামাজের \textbf{\lat{difficulty} (কঠিনতা)} খুশুর অভাবে; যার অন্তর আল্লাহর সামনে বিনয়ী, তার নামাজ সহজ ও উপকারী। (আস-সা'দী, সূরা বাকারাহ)
\end{itemize}
\subsubsection*{৭. সংশ্লিষ্ট হাদিস}

\begin{itemize}
\item \textbf{হাদিস (আবু দাউদ ১৩১৯):} নবীজি (সা.) বিপদে পড়লে নামাজে দাঁড়িয়ে যেতেন। (সনদ নিয়ে আলেমদের আলোচনা আছে; ইমাম আহমাদ ও হাসান বাসরীর বক্তব্যেও এ কথা আছে।)
\item \textbf{আয়াত (বাকারাহ ২:১৫৩):} \lat{“}হে ঈমানদারগণ! তোমরা ধৈর্য ও নামাজের মাধ্যমে সাহায্য চাও...\lat{”}
\end{itemize}
\subsubsection*{৮. গভীর শিক্ষা (\lat{Deep} \lat{Lessons})}

\begin{enumerate}
\item \textbf{নামাজ — বিপদের ওষুধ।} দুশ্চিন্তা, অস্থিরতা, বিপদ — এসবের মুখে নামাজ আল্লাহর কাছে সাহায্যের দরজা।
\item \textbf{খুশু নামাজকে সহজ করে।} নামাজ কঠিন লাগে তখনই, যখন খুশু নেই; খুশু নামাজের \textbf{\lat{sweetness} (মাধুর্য)}।
\item \textbf{সবর ও নামাজ জোড়া।} ধৈর্য ও নামাজ — দুইটি মিলে মুমিনের শক্তি।
\end{enumerate}
\subsubsection*{৯. জীবনে প্রয়োগ (\lat{Practical} \lat{Application})}

\begin{enumerate}
\item কঠিন সময়ে নামাজ ও দোয়াকে প্রথম আশ্রয় বানান — হতাশা নয়।
\item নামাজে খুশু বাড়ানোর চর্চা করুন — অর্থ বুঝে পড়ুন, ধীরে পড়ুন।
\end{enumerate}
\medskip\hrule\medskip

\subsection*{আয়াত ৫ — সূরা আল-আনকাবুত (২৯) — আয়াত ৪৫}

\subsubsection*{১. সূত্র কার্ড}

\begin{table}[h]\centering\small
\begin{tabular}{ll}
বিষয় & বিবরণ \\
\hline
\textbf{সূরা} & আল-আনকাবুত (২৯) — মক্কী \\
\textbf{আয়াত} & ৪৫ \\
\textbf{পারা} & ২১ \\
\textbf{মূল বিষয়} & নামাজ অশ্লীলতা ও মন্দ কাজ থেকে বিরত রাখে \\
\textbf{বক্তা} & আল্লাহ — নবীজিকে সম্বোধন \\
\hline
\end{tabular}
\end{table}

\subsubsection*{২. প্রসঙ্গ ও প্রেক্ষাপট (\lat{Context})}

আয়াতটি নবীজিকে নির্দেশ দিচ্ছে — কিতাব তিলাওয়াত করো ও নামাজ কায়েম করো। তারপর নামাজের সবচেয়ে বিখ্যাত ফলাফল ঘোষিত হয়েছে: \lat{“}নিশ্চয়ই নামাজ অশ্লীল ও মন্দ কাজ থেকে বিরত রাখে।\lat{”} অর্থাৎ নামাজ কেবল ইবাদত নয় — এটি একটি নৈতিক \textbf{\lat{shield} (ঢাল)}; যিনি সঠিকভাবে নামাজ পড়েন, তার জীবন থেকে ফাহেশা ও মুনকার দূর হয়।

\subsubsection*{৩. মূল আরবি}

\begin{center}\arab{اتْلُ مَا أُوحِيَ إِلَيْكَ مِنَ الْكِتَابِ وَأَقِمِ الصَّلَاةَ ۖ إِنَّ الصَّلَاةَ تَنْهَىٰ عَنِ الْفَحْشَاءِ وَالْمُنكَرِ ۗ وَلَذِكْرُ اللَّهِ أَكْبَرُ ۗ وَاللَّهُ يَعْلَمُ مَا تَصْنَعُونَ}\end{center}

\textbf{উচ্চারণ:} উৎলু মা- উহিয়া ইলাইকা মিনাল কিতা-বি ওয়া আক্বিমিস-সালা-হ, ইন্নাস-সালা-তা তানহা- 'আনিল ফাহশা-য়ি ওয়াল মুনকার, ওয়া লাযিকরুল্লা-হি আকবার, ওয়াল্লা-হু ইয়া'লামু মা- তাসনা'ঊন।

\subsubsection*{৪. বাংলা অর্থ (\lat{pure})}

\textbf{\lat{“}তুমি কিতাব থেকে যা ওহী করা হয়েছে তা তিলাওয়াত করো এবং নামাজ কায়েম করো। নিশ্চয়ই নামাজ অশ্লীল ও মন্দ কাজ থেকে বিরত রাখে। আর আল্লাহর স্মরণই সর্বশ্রেষ্ঠ। আল্লাহ জানেন যা তোমরা করো।\lat{”}}

\subsubsection*{৫. শব্দের ব্যাখ্যা (\lat{Vocabulary})}

\begin{table}[h]\centering\small
\begin{tabular}{ll}
শব্দ & অর্থ \\
\hline
\textbf{\arab{اتْلُ}} & তিলাওয়াত করো — পাঠ করো \\
\textbf{\arab{الْفَحْشَاءِ}} & অশ্লীলতা — চরম অন্যায়, ব্যভিচার ইত্যাদি \\
\textbf{\arab{الْمُنكَرِ}} & মন্দ কাজ — শরীয়ত-অস্বীকৃত কাজ \\
\textbf{\arab{تَنْهَىٰ}} & বিরত রাখে — নিষেধ করে, বাধা দেয় \\
\textbf{\arab{الذِّكْرِ}} & স্মরণ — আল্লাহর যিকর \\
\hline
\end{tabular}
\end{table}

\subsubsection*{৬. গভীর তাফসীর (\lat{Deep} \lat{Tafsir})}

\begin{itemize}
\item \textbf{ইবনে কাসীর:} নামাজ যদি সঠিকভাবে আদায় হয়, তবে তা ফাহেশা ও মুনকার থেকে বিরত রাখে। ইবনে আব্বাস (রাঃ) বলেন — নামাজে খুশু ও অন্তরের উপস্থিতি না থাকলে নামাজ তাকে বিরত রাখে না। ইবনে মাসউদ (রাঃ)-এর ভাষায় — যে নামাজ তাকে অন্যায় থেকে বিরত রাখে না, তার নামাজ কেবল \textbf{\lat{delay} (বিলম্ব)}। (তাফসীর ইবনে কাসীর, সূরা আনকাবুত)
\item \textbf{কুরতুবী:} \lat{“}\arab{وَلَذِكْرُ} \arab{اللَّهِ} \arab{أَكْبَرُ}\lat{”} — আল্লাহর স্মরণ সবচেয়ে বড়; কেউ কেউ বলেন অর্থ — নামাজের মধ্যে আল্লাহর যিকরই সর্বশ্রেষ্ঠ অংশ; কারো মতে আল্লাহ তোমাদের স্মরণ করাই সবচেয়ে বড়। (তাফসীরুল কুরতুবী, সূরা আনকাবুত)
\item \textbf{সা'দী:} নামাজের প্রকৃত ফল — অন্তরে আল্লাহর ভয় ও লজ্জা জন্মে, যা মানুষকে পাপ থেকে বিরত রাখে। নামাজ যত খুশুপূর্ণ, প্রতিরোধ ক্ষমতা তত শক্তিশালী। (আস-সা'দী, সূরা আনকাবুত)
\end{itemize}
\subsubsection*{৭. সংশ্লিষ্ট হাদিস}

\begin{itemize}
\item \textbf{হাদিস (আবু দাউদ ৭৯৬, হাসান):} \lat{“}মানুষ নামাজ পড়ে, কিন্তু তার নামাজের থেকে শুধু দশমাংশ, নবমাংশ... লেখা হয়।\lat{”} — নামাজের প্রভাব নামাজের মানের উপর নির্ভরশীল।
\item \textbf{আয়াত (মু'মিনুন ২৩:২):} \lat{“}যারা নিজেদের নামাজে বিনয়ী\lat{”} — খুশুর সাথে নামাজের সম্পর্ক।
\end{itemize}
\subsubsection*{৮. গভীর শিক্ষা (\lat{Deep} \lat{Lessons})}

\begin{enumerate}
\item \textbf{নামাজ — নৈতিক বাধা।} নামাজ যত সঠিক, পাপের প্রতি প্রবণতা তত কম — এ আল্লাহর প্রতিশ্রুতি।
\item \textbf{নামাজের প্রভাব নামাজের মানের উপর।} যে নামাজ খুশুহীন, সে নামাজ পাপ থেকে বিরত রাখার শক্তি হারায়।
\item \textbf{আল্লাহর স্মরণ শ্রেষ্ঠ।} নামাজের ভেতরে ও বাইরে আল্লাহর যিকরই সবচেয়ে বড় ইবাদত।
\end{enumerate}
\subsubsection*{৯. জীবনে প্রয়োগ (\lat{Practical} \lat{Application})}

\begin{enumerate}
\item নামাজের মান যাচাই করুন — নামাজের পরে কোনো পাপে \textbf{\lat{attracted} (আকৃষ্ট)} হলে নামাজের মান নিয়ে ভাবুন।
\item দিনের শুরুতে নামাজকে বর্ম বানান — নামাজই পাপের বিরুদ্ধে আপনার প্রতিরোধ।
\end{enumerate}
\medskip\hrule\medskip

\subsection*{আয়াত ৬ — সূরা আল-মু'মিনুন (২৩) — আয়াত ১-২}

\subsubsection*{১. সূত্র কার্ড}

\begin{table}[h]\centering\small
\begin{tabular}{ll}
বিষয় & বিবরণ \\
\hline
\textbf{সূরা} & আল-মু'মিনুন (২৩) — মক্কী \\
\textbf{আয়াত} & ১-২ \\
\textbf{পারা} & ১৮ \\
\textbf{মূল বিষয়} & সফল মুমিনের প্রথম গুণ — নামাজে খুশু \\
\textbf{বক্তা} & আল্লাহ \\
\hline
\end{tabular}
\end{table}

\subsubsection*{২. প্রসঙ্গ ও প্রেক্ষাপট (\lat{Context})}

সূরার শুরুতে মুমিনদের সাফল্যের ঘোষণা — \lat{“}নিশ্চয়ই সফল হয়েছে মুমিনরা\lat{”} — তারপর তাদের গুণাবলি গণনা করা হয়েছে; প্রথম গুণই হলো নামাজে খুশু। অর্থাৎ আল্লাহর কাছে সফলতার তালিকায় নামাজের খুশু সবচেয়ে আগে। আয়াত ৯-এ আবার নামাজের \textbf{যত্ন (\lat{guardianship})} — \lat{“}যারা নিজেদের নামাজের প্রতি যত্নবান\lat{”} — সাফল্যের গুণ হিসেবে উল্লেখ আছে।

\subsubsection*{৩. মূল আরবি}

\begin{center}\arab{قَدْ أَفْلَحَ الْمُؤْمِنُونَ ۖ الَّذِينَ هُمْ فِي صَلَاتِهِمْ خَاشِعُونَ}\end{center}

\textbf{উচ্চারণ:} ক্বাদ আফলাহাল মু'মিনূন। আল্লাযীনা হুম ফী সালা-তিহিম খা-শি'ঊন।

\subsubsection*{৪. বাংলা অর্থ (\lat{pure})}

\textbf{\lat{“}নিশ্চয়ই সফল হয়েছে মুমিনরা — যারা নিজেদের নামাজে বিনয়ী (খুশুকারী)।\lat{”}}

\subsubsection*{৫. শব্দের ব্যাখ্যা (\lat{Vocabulary})}

\begin{table}[h]\centering\small
\begin{tabular}{ll}
শব্দ & অর্থ \\
\hline
\textbf{\arab{أَفْلَحَ}} & সফল হয়েছে — সাফল্য লাভ করেছে \\
\textbf{\arab{الْخَاشِعُونَ}} & বিনয়ীরা — যাদের অন্তর ও অঙ্গ-প্রত্যঙ্গ আল্লাহর সামনে নত \\
\hline
\end{tabular}
\end{table}

\subsubsection*{৬. গভীর তাফসীর (\lat{Deep} \lat{Tafsir})}

\begin{itemize}
\item \textbf{ইবনে কাসীর:} আলী (রাঃ) ও ইবনে আব্বাস (রাঃ) বলেন — খুশু মানে নামাজে অন্তরের উপস্থিতি; অঙ্গ-প্রত্যঙ্গ স্থির, দৃষ্টি সিজদার দিকে। নামাজে খুশু সফলতার প্রথম চাবি। (তাফসীর ইবনে কাসীর, সূরা মু'মিনুন)
\item \textbf{তাবারী:} খুশু দুই ধরনের — অন্তরের খুশু (আল্লাহর ভয় ও বিনয়) এবং অঙ্গের খুশু (স্থিরতা, না তাকানো)। নামাজে উভয়ই কাম্য। (জামিউল বায়ান, সূরা মু'মিনুন)
\item \textbf{সা'দী:} খুশু হলো নামাজের \textbf{\lat{soul} (প্রাণ)}; খুশুহীন নামাজ শরীরের ব্যায়ামে পরিণত হয়। খুশু আনার উপায় — নামাজে যা পড়ি তা বোঝা এবং আল্লাহর সামনে দাঁড়ানোর অনুভূতি। (আস-সা'দী, সূরা মু'মিনুন)
\end{itemize}
\subsubsection*{৭. সংশ্লিষ্ট হাদিস}

\begin{itemize}
\item \textbf{হাদিস (মুসলিম ৪৮২):} \lat{“}বান্দা তার রবের সবচেয়ে নিকটবর্তী সিজদার সময় — তাই সিজদায় দোয়া বেশি করো।\lat{”} — সিজদার খুশুই নৈকট্যের কারণ।
\item \textbf{হাদিস (বুখারি ৭৫৭, মুসলিম ৩৯৭):} \lat{“}ফিরে যাও, নামাজ পড়ো — তুমি নামাজ পড়োনি\lat{”} (মুসি' সালাতুহু হাদিস) — খুশু ও ইতমিনানহীন নামাজ আল্লাহর কাছে গণ্য নয়।
\end{itemize}
\subsubsection*{৮. গভীর শিক্ষা (\lat{Deep} \lat{Lessons})}

\begin{enumerate}
\item \textbf{খুশু — সাফল্যের প্রথম শর্ত।} আল্লাহর তালিকায় খুশু প্রথমে — অর্থাৎ নামাজের \textbf{\lat{quality} (মান)} সংখ্যার চেয়ে গুরুত্বপূর্ণ।
\item \textbf{খুশু অর্জনযোগ্য।} খুশু জন্মগত নয় — অভ্যাস ও চর্চায় আসে; অর্থবোঝা, ধীরতা, দৃষ্টি নিয়ন্ত্রণ।
\item \textbf{অন্তর ও অঙ্গের মিল।} অন্তরের বিনয় অঙ্গের স্থিরতায় প্রকাশ পায়।
\end{enumerate}
\subsubsection*{৯. জীবনে প্রয়োগ (\lat{Practical} \lat{Application})}

\begin{enumerate}
\item নামাজে দৃষ্টি সিজদার স্থানে রাখুন — চোখ এদিক-ওদিক না ঘুরান।
\item নামাজের অর্থ বুঝে পড়ার চেষ্টা করুন — অন্তত সূরা ফাতিহা ও সাধারণ সূরাগুলোর অর্থ।
\end{enumerate}
\medskip\hrule\medskip

\subsection*{আয়াত ৭ — সূরা আল-বাকারাহ (২) — আয়াত ২৩৮}

\subsubsection*{১. সূত্র কার্ড}

\begin{table}[h]\centering\small
\begin{tabular}{ll}
বিষয় & বিবরণ \\
\hline
\textbf{সূরা} & আল-বাকারাহ (২) — মাদানী \\
\textbf{আয়াত} & ২৩৮ \\
\textbf{পারা} & ২ \\
\textbf{মূল বিষয়} & সব নামাজের প্রতি যত্ন — বিশেষ করে আসর (মধ্যবর্তী নামাজ) \\
\textbf{বক্তা} & আল্লাহ \\
\hline
\end{tabular}
\end{table}

\subsubsection*{২. প্রসঙ্গ ও প্রেক্ষাপট (\lat{Context})}

আয়াতটি মুমিনদের সব নামাজের প্রতি যত্নবান হওয়ার নির্দেশ দিচ্ছে, বিশেষভাবে \lat{“}মধ্যবর্তী নামাজ\lat{”} (সালাতুল উস্তা)-এর কথা বলেছে। অধিকাংশ আলেমের মতে এটি \textbf{আসরের নামাজ} (বুখারি ৫৫৪-এর হাদিস অনুযায়ী)। আসরের গুরুত্ব এত বেশি যে, নামাজ ত্যাগকারীর ভয়াবহ পরিণতিও আসরের সাথে যুক্ত (বুখারি ৫৫২)।

\subsubsection*{৩. মূল আরবি}

\begin{center}\arab{حَافِظُوا عَلَى الصَّلَوَاتِ وَالصَّلَاةِ الْوُسْطَىٰ وَقُومُوا لِلَّهِ قَانِتِينَ}\end{center}

\textbf{উচ্চারণ:} হা-ফিযূ 'আলাস-সালাওয়া-তি ওয়াস-সালা-তিল উসত্বা, ওয়া কূমূলিল্লা-হি কা-নিতীন।

\subsubsection*{৪. বাংলা অর্থ (\lat{pure})}

\textbf{\lat{“}তোমরা সব নামাজের প্রতি যত্নবান হও, বিশেষ করে মধ্যবর্তী নামাজের (আসর)। আর আল্লাহর জন্য বিনীতভাবে দাঁড়াও।\lat{”}}

\subsubsection*{৫. শব্দের ব্যাখ্যা (\lat{Vocabulary})}

\begin{table}[h]\centering\small
\begin{tabular}{ll}
শব্দ & অর্থ \\
\hline
\textbf{\arab{حَافِظُوا}} & যত্নবান হও — হিফাজত করো \\
\textbf{\arab{الصَّلَوَاتِ}} & নামাজসমূহ — পাঁচ ওয়াক্ত \\
\textbf{\arab{الْوُسْطَىٰ}} & মধ্যবর্তী — মাঝখানের (আসর) \\
\textbf{\arab{قَانِتِينَ}} & বিনীত — আল্লাহর সামনে নত, স্থির \\
\hline
\end{tabular}
\end{table}

\subsubsection*{৬. গভীর তাফসীর (\lat{Deep} \lat{Tafsir})}

\begin{itemize}
\item \textbf{ইবনে কাসীর:} \lat{“}সালাতুল উস্তা\lat{”} সম্পর্কে সাহাবাদের মতভেদ আছে — অধিকাংশ (আলী, ইবনে আব্বাস, আবু হুরাইরা, উমার প্রমুখ) বলেছেন আসর। নবীজির হাদিসে — \lat{“}আমরা আসরের নামাজকে 'উস্তা' বলতাম\lat{”} (বুখারি ৫৫৪)। (তাফসীর ইবনে কাসীর, সূরা বাকারাহ)
\item \textbf{কুরতুবী:} আয়াতে নামাজের হিফাজতের নির্দেশ — সময়মতো, পরিপূর্ণভাবে, জামাতে। \lat{“}কুনুত\lat{”} অর্থ বিনয়ের সাথে দাঁড়ানো — নামাজে নীরব স্থিরতা। (তাফসীরুল কুরতুবী, সূরা বাকারাহ)
\item \textbf{মাআরিফুল কুরআন:} আসরের নামাজ কর্মব্যস্ত দিনের শেষে আসে — এ সময় মানুষ অলস ও ব্যস্ত থাকে; তাই বিশেষ তাকিদ। (মাআরিফুল কুরআন, সূরা বাকারাহ)
\end{itemize}
\subsubsection*{৭. সংশ্লিষ্ট হাদিস}

\begin{itemize}
\item \textbf{হাদিস (বুখারি ৫৫২, মুসলিম ৬২৬):} \lat{“}যে ব্যক্তি আসরের নামাজ ছেড়ে দিল, সে যেন পরিবার ও ধন-সম্পদ হারিয়ে ফেলল।\lat{”}
\item \textbf{হাদিস (বুখারি ৫৫৪):} নবীজি (সা.) বলেছেন — \lat{“}আমরা সালাতুল উস্তা বলতাম আসরকে।\lat{”}
\end{itemize}
\subsubsection*{৮. গভীর শিক্ষা (\lat{Deep} \lat{Lessons})}

\begin{enumerate}
\item \textbf{সব নামাজ সমান গুরুত্বপূর্ণ।} আয়াত \lat{“}সব নামাজের প্রতি\lat{”} দিয়ে শুরু — কোনো নামাজ ছোট নয়।
\item \textbf{আসরের বিশেষত্ব।} কর্মব্যস্ত দিনের শেষে আসর — এই সময় নামাজ ছুটে যাওয়ার ঝুঁকি বেশি; তাই বিশেষ তাকিদ।
\item \textbf{নামাজে বিনয়।} \lat{“}\arab{قَانِتِينَ}\lat{”} — নামাজে দাঁড়ানো মানে আল্লাহর সামনে নত হওয়া।
\end{enumerate}
\subsubsection*{৯. জীবনে প্রয়োগ (\lat{Practical} \lat{Application})}

\begin{enumerate}
\item আসরের নামাজের জন্য দিনের মধ্যে \textbf{\lat{reminder} (স্মরণক)} রাখুন — ব্যস্ততায় নামাজ যেন না ছুটে।
\item সব নামাজ সমান যত্নে পড়ুন — একটিকে বাদ দিয়ে অন্যটি নয়।
\end{enumerate}
\medskip\hrule\medskip

\subsection*{আয়াত ৮ — সূরা হুদ (১১) — আয়াত ১১৪}

\subsubsection*{১. সূত্র কার্ড}

\begin{table}[h]\centering\small
\begin{tabular}{ll}
বিষয় & বিবরণ \\
\hline
\textbf{সূরা} & হুদ (১১) — মক্কী \\
\textbf{আয়াত} & ১১৪ \\
\textbf{পারা} & ১২ \\
\textbf{মূল বিষয়} & দিনের দুই প্রান্তে নামাজ; নেকি মন্দকে দূর করে \\
\textbf{বক্তা} & আল্লাহ — নবীজিকে সম্বোধন \\
\hline
\end{tabular}
\end{table}

\subsubsection*{২. প্রসঙ্গ ও প্রেক্ষাপট (\lat{Context})}

নবীজিকে নির্দেশ — দিনের দুই প্রান্তে (ফজর, মাগরিব) ও রাতের প্রথমাংশে (ইশা) নামাজ কায়েম করো। তারপর গুরুত্বপূর্ণ শিক্ষা: \lat{“}নিশ্চয়ই নেক আমল মন্দ আমলকে দূর করে দেয়।\lat{”} আলেমরা বলেন — এটি তাওবার জন্য সাধারণ নির্দেশ; নামাজসহ নেক আমল পূর্বের গুনাহ মিটিয়ে দেয়। নামাজ এখানে গুনাহ মোচনের \textbf{\lat{means} (উপায়)} হিসেবে বর্ণিত।

\subsubsection*{৩. মূল আরবি}

\begin{center}\arab{وَأَقِمِ الصَّلَاةَ طَرَفَيِ النَّهَارِ وَزُلَفًا مِّنَ اللَّيْلِ ۚ إِنَّ الْحَسَنَاتِ يُذْهِبْنَ السَّيِّئَاتِ ۚ ذَٰلِكَ ذِكْرَىٰ لِلذَّاكِرِينَ}\end{center}

\textbf{উচ্চারণ:} ওয়া আক্বিমিস-সালা-তা ত্বারাফাইন্নাহা-রি ওয়া যুলাফাম মিনাল্লাইল, ইন্নাল হাসানা-তি ইউযহিবনাস সাইয়ি-আ-তি, যা-লিকা যিকরা- লিয-যা-কিরীন।

\subsubsection*{৪. বাংলা অর্থ (\lat{pure})}

\textbf{\lat{“}আর তুমি দিনের দুই প্রান্তে এবং রাতের কিছু অংশে নামাজ কায়েম করো। নিশ্চয়ই নেক কাজগুলো মন্দ কাজগুলোকে দূর করে দেয়। এটি স্মরণকারীদের জন্য উপদেশ।\lat{”}}

\subsubsection*{৫. শব্দের ব্যাখ্যা (\lat{Vocabulary})}

\begin{table}[h]\centering\small
\begin{tabular}{ll}
শব্দ & অর্থ \\
\hline
\textbf{\arab{طَرَفَيِ} \arab{النَّهَارِ}} & দিনের দুই প্রান্ত — ভোর ও সন্ধ্যা \\
\textbf{\arab{زُلَفًا} \arab{مِّنَ} \arab{اللَّيْلِ}} & রাতের কিছু অংশ — রাতের প্রথমাংশ (ইশা) \\
\textbf{\arab{الْحَسَنَاتِ}} & নেক কাজ — সৎকর্ম \\
\textbf{\arab{يُذْهِبْنَ}} & দূর করে — মুছে দেয় \\
\textbf{\arab{السَّيِّئَاتِ}} & মন্দ কাজ — পাপ \\
\hline
\end{tabular}
\end{table}

\subsubsection*{৬. গভীর তাফসীর (\lat{Deep} \lat{Tafsir})}

\begin{itemize}
\item \textbf{ইবনে কাসীর:} ইবনে আব্বাস (রাঃ) বলেন — \lat{“}দিনের দুই প্রান্ত\lat{”} অর্থ ফজর ও মাগরিব; \lat{“}রাতের কিছু অংশ\lat{”} অর্থ ইশা। হাদিসে এসেছে — নবীজি বলেছেন: এই আয়াতের সাথে কোনো হাদিসের তুলনা নেই (অর্থাৎ নেকি-মন্দ দূরীকরণের এই বাণী সবচেয়ে সুন্দর) — বর্ণনাটি তিরমিযীতে আছে (হাসান)। (তাফসীর ইবনে কাসীর, সূরা হুদ)
\item \textbf{তাবারী:} নামাজের পর নামাজ, নেক আমলের পর নেক আমল — পূর্বের মন্দকে দূর করে; এটি তাওবার পথ খুলে দেয়। (জামিউল বায়ান, সূরা হুদ)
\item \textbf{সা'দী:} আয়াতের শিক্ষা — নামাজসহ নেক আমল পাপ মুছে দেয়; বিশেষত পাপের পর সাথে সাথে নেক আমল করলে তা পাপের প্রভাব দূর করে। (আস-সা'দী, সূরা হুদ)
\end{itemize}
\subsubsection*{৭. সংশ্লিষ্ট হাদিস}

\begin{itemize}
\item \textbf{হাদিস (মুসলিম ২৩৩):} \lat{“}পাঁচ ওয়াক্ত নামাজ এবং জুমা থেকে জুমা — মধ্যবর্তী গুনাহের কাফফারা, যতক্ষণ না বড় গুনাহ করা হয়।\lat{”}
\item \textbf{হাদিস (তিরমিযী ৩১৪৭, হাসান):} নবীজি বলেছেন — \lat{“}তুমি যেখানেই থাকো, আল্লাহকে ভয় করো; পাপের পর নেক আমল করো, তা পাপ মুছে দেবে।\lat{”}
\end{itemize}
\subsubsection*{৮. গভীর শিক্ষা (\lat{Deep} \lat{Lessons})}

\begin{enumerate}
\item \textbf{নামাজ গুনাহ মোচনের মাধ্যম।} আয়াতের প্রতিশ্রুতি — নেক আমল (বিশেষত নামাজ) মন্দ আমলকে সরিয়ে দেয়।
\item \textbf{দিন-রাতের কাঠামো।} ফজর, মাগরিব, ইশা — দিনের শুরু, শেষ ও রাতের প্রহর; নামাজ সারাদিনকে ঢেকে রাখে।
\item \textbf{হতাশা নয়, সংশোধন।} পাপের পর হতাশ না হয়ে নেক আমলে ফেরা — আল্লাহর রহমতের দরজা।
\end{enumerate}
\subsubsection*{৯. জীবনে প্রয়োগ (\lat{Practical} \lat{Application})}

\begin{enumerate}
\item পাপ হয়ে গেলে দেরি না করে নামাজ, তাওবা ও নেক আমলে \textbf{\lat{return} (ফিরে আসুন)}।
\item ফজর ও মাগরিবের নামাজকে দিনের অপরিহার্য অংশ করুন।
\end{enumerate}
\medskip\hrule\medskip

\subsection*{আয়াত ৯ — সূরা বনী ইসরাঈল / আল-ইসরা (১৭) — আয়াত ৭৮-৭৯}

\subsubsection*{১. সূত্র কার্ড}

\begin{table}[h]\centering\small
\begin{tabular}{ll}
বিষয় & বিবরণ \\
\hline
\textbf{সূরা} & বনী ইসরাঈল / আল-ইসরা (১৭) — মক্কী \\
\textbf{আয়াত} & ৭৮-৭৯ \\
\textbf{পারা} & ১৫ \\
\textbf{মূল বিষয়} & নামাজের সময়সীমা; ফজরের কুরআন; রাতের তাহাজ্জুদ \\
\textbf{বক্তা} & আল্লাহ — নবীজিকে সম্বোধন \\
\hline
\end{tabular}
\end{table}

\subsubsection*{২. প্রসঙ্গ ও প্রেক্ষাপট (\lat{Context})}

আয়াত ৭৮-এ নামাজের সময়সীমা বলা হয়েছে — \lat{“}সূর্য ঢলার সময় (যুহর-আসর) থেকে রাতের অন্ধকার পর্যন্ত (মাগরিব-ইশা)\lat{”} এবং ফজরের কুরআন (ফজরের নামাজ), যা ফেরেশতারা সাক্ষী থাকে। আয়াত ৭৯-এ নবীজিকে রাতের তাহাজ্জুদের নির্দেশ — \lat{“}রাতের কিছু অংশে জাগ্রত থেকে তাহাজ্জুদ পড়ো\lat{”} — যা নবীজির জন্য অতিরিক্ত (নাফিল) হিসেবে বিশেষ মর্যাদা এনে দেয় (\lat{“}মাকামে মাহমুদ\lat{”})।

\subsubsection*{৩. মূল আরবি}

\begin{center}\arab{أَقِمِ الصَّلَاةَ لِدُلُوكِ الشَّمْسِ إِلَىٰ غَسَقِ اللَّيْلِ وَقُرْآنَ الْفَجْرِ ۖ إِنَّ قُرْآنَ الْفَجْرِ كَانَ مَشْهُودًا وَمِنَ اللَّيْلِ فَتَهَجَّدْ بِهِ نَافِلَةً لَّكَ ۖ عَسَىٰ أَن يَبْعَثَكَ رَبُّكَ مَقَامًا مَّحْمُودًا}\end{center}

\textbf{উচ্চারণ:} আক্বিমিস-সালা-তা লিদুলূকিশ শামসি ইলা- গাসাকিল্লাইলি ওয়া কুরআ-নাল ফাজর, ইন্না কুরআ-নাল ফাজরি কা-না মাশহূদা। ওয়া মিনাল্লাইলি ফাতাহাজ্জাদ বিহী না-ফিলাতাল্লাকা, 'আসা- আই ইয়াব'আসাকা রাব্বুকা মাকা-মাম মাহমূদা।

\subsubsection*{৪. বাংলা অর্থ (\lat{pure})}

\textbf{\lat{“}তুমি সূর্য ঢলে পড়ার সময় থেকে রাতের অন্ধকার পর্যন্ত নামাজ কায়েম করো, আর (কায়েম করো) ফজরের কুরআন (নামাজ) — নিশ্চয়ই ফজরের কুরআন (নামাজ) সাক্ষ্যপ্রাপ্ত। আর রাতের কিছু অংশে জাগ্রত থেকে তাহাজ্জুদ পড়ো — এটি তোমার জন্য অতিরিক্ত (নাফিল); আশা করা যায় তোমার রব তোমাকে মর্যাদাপূর্ণ স্থানে (মাকামে মাহমুদ) পৌঁছাবেন।\lat{”}}

\subsubsection*{৫. শব্দের ব্যাখ্যা (\lat{Vocabulary})}

\begin{table}[h]\centering\small
\begin{tabular}{ll}
শব্দ & অর্থ \\
\hline
\textbf{\arab{دُلُوكِ} \arab{الشَّمْسِ}} & সূর্যের ঢল — মধ্যাহ্নের পর \\
\textbf{\arab{غَسَقِ} \arab{اللَّيْلِ}} & রাতের অন্ধকার — রাতের ঘনত্ব \\
\textbf{\arab{قُرْآنَ} \arab{الْفَجْرِ}} & ফজরের কুরআন — ফজরের নামাজ \\
\textbf{\arab{مَشْهُودًا}} & সাক্ষ্যপ্রাপ্ত — দিবা-রাত্রির ফেরেশতারা উপস্থিত থাকে \\
\textbf{\arab{فَتَهَجَّدْ}} & তাহাজ্জুদ পড়ো — রাত জেগে নামাজ \\
\textbf{\arab{مَقَامًا} \arab{مَّحْمُودًا}} & প্রশংসিত স্থান — মাকামে মাহমুদ \\
\hline
\end{tabular}
\end{table}

\subsubsection*{৬. গভীর তাফসীর (\lat{Deep} \lat{Tafsir})}

\begin{itemize}
\item \textbf{ইবনে কাসীর:} \lat{“}দুলুক\lat{”} মানে সূর্যের ঢল — যুহরের সময় শুরু; থেকে \lat{“}গাসাক\lat{”} মানে রাতের অন্ধকার — ইশার সময় পর্যন্ত। ফজরের নামাজের বিশেষত্ব — দিবা ও রাত্রির ফেরেশতারা তাতে উপস্থিত থাকে। (তাফসীর ইবনে কাসীর, সূরা ইসরা)
\item \textbf{কুরতুবী:} আয়াত ৭৯-এ তাহাজ্জুদের নির্দেশ — নবীজির উপর এটা ওয়াজিব ছিল (কেউ কেউ বলেন), উম্মতের জন্য সুন্নত। \lat{“}মাকামে মাহমুদ\lat{”} — কিয়ামতের দিনের মহান সুপারিশের স্থান। (তাফসীরুল কুরতুবী, সূরা ইসরা)
\item \textbf{মাআরিফুল কুরআন:} এই আয়াত থেকেই তাহাজ্জুদ নামাজের বিধান ও মাকামে মাহমুদের প্রতিশ্রুতি — রাতের নামাজ মুমিনের বিশেষ মর্যাদার কারণ। (মাআরিফুল কুরআন, সূরা ইসরা)
\end{itemize}
\subsubsection*{৭. সংশ্লিষ্ট হাদিস}

\begin{itemize}
\item \textbf{হাদিস (মুসলিম ১১৬৩):} \lat{“}ফরজ নামাজের পর সর্বোত্তম নামাজ হলো রাতের নামাজ (তাহাজ্জুদ)।\lat{”}
\item \textbf{হাদিস (বুখারি ৪৮৩৬, মুসলিম ২৮১৯):} নবীজি রাতের নামাজে এত দাঁড়াতেন যে পা ফুলে যেত; জিজ্ঞেস করলে বললেন — \lat{“}আমি কি কৃতজ্ঞ বান্দা হব না?\lat{”}
\end{itemize}
\subsubsection*{৮. গভীর শিক্ষা (\lat{Deep} \lat{Lessons})}

\begin{enumerate}
\item \textbf{নামাজের সময় কাঠামো।} আয়াতটি যুহর থেকে ইশা পর্যন্ত সময়সীমা দিয়েছে — সময়ের প্রতি যত্ন।
\item \textbf{ফজরের বিশেষত্ব।} ফজরের নামাজ ফেরেশতাদের সাক্ষ্যপ্রাপ্ত — ভোরের ইবাদতের মর্যাদা।
\item \textbf{তাহাজ্জুদের মর্যাদা।} রাতের নামাজ নবীজির বিশেষ আমল ও মাকামে মাহমুদের পথ — মুমিনের জন্যও সুন্নত।
\end{enumerate}
\subsubsection*{৯. জীবনে প্রয়োগ (\lat{Practical} \lat{Application})}

\begin{enumerate}
\item ফজরের নামাজ জামাতে পড়ার চেষ্টা করুন — ফেরেশতাদের সাক্ষ্যের সুবর্ণ সুযোগ।
\item সপ্তাহে ২-৩ রাত অন্তত ২ রাকাত তাহাজ্জুদ পড়ার \textbf{\lat{habit} (অভ্যাস)} গড়ুন।
\end{enumerate}
\medskip\hrule\medskip

\subsection*{আয়াত ১০ — সূরা আল-মাউন (১০৭) — আয়াত ৪-৫}

\subsubsection*{১. সূত্র কার্ড}

\begin{table}[h]\centering\small
\begin{tabular}{ll}
বিষয় & বিবরণ \\
\hline
\textbf{সূরা} & আল-মাউন (১০৭) — মক্কী \\
\textbf{আয়াত} & ৪-৫ \\
\textbf{পারা} & ৩০ \\
\textbf{মূল বিষয়} & নামাজে উদাসীনদের ধিক্কার — সতর্কবার্তা \\
\textbf{বক্তা} & আল্লাহ \\
\hline
\end{tabular}
\end{table}

\subsubsection*{২. প্রসঙ্গ ও প্রেক্ষাপট (\lat{Context})}

সূরায় এমন লোকদের নিন্দা করা হয়েছে যারা দ্বীনের দিন (হিসাব) অস্বীকার করে, এতিমকে ধাক্কা দেয় এবং মিসকিনকে খাওয়ায় না। আয়াত ৪-৫-এ বলা হয়েছে — \lat{“}ধিক সেই নামাজিদের জন্য, যারা নিজেদের নামাজ সম্পর্কে উদাসীন।\lat{”} আলেমদের মতে \lat{“}সাহূন\lat{”} অর্থ নামাজের সময় ভুলে যাওয়া বা নামাজে অবহেলা — এখানে এমন নামাজিদের ভয়াবহ সতর্কবার্তা, যারা নামাজের গুরুত্ব বোঝে না।

\subsubsection*{৩. মূল আরবি}

\begin{center}\arab{فَوَيْلٌ لِّلْمُصَلِّينَ ۖ الَّذِينَ هُمْ عَن صَلَاتِهِمْ سَاهُونَ}\end{center}

\textbf{উচ্চারণ:} ফাওয়াইলুল লিল-মুসাল্লীন। আল্লাযীনা হুম 'আন সালা-তিহিম সা-হূন।

\subsubsection*{৪. বাংলা অর্থ (\lat{pure})}

\textbf{\lat{“}অতএব ধিক সেই নামাজিদের জন্য — যারা নিজেদের নামাজ সম্পর্কে উদাসীন।\lat{”}}

\subsubsection*{৫. শব্দের ব্যাখ্যা (\lat{Vocabulary})}

\begin{table}[h]\centering\small
\begin{tabular}{ll}
শব্দ & অর্থ \\
\hline
\textbf{\arab{وَيْلٌ}} & ধিক্ — ধ্বংস/দুর্ভোগ \\
\textbf{\arab{لِلْمُصَلِّينَ}} & নামাজিদের জন্য — যারা নামাজ পড়ে \\
\textbf{\arab{سَاهُونَ}} & উদাসীন — অবহেলা করে, ভুলে যায় \\
\hline
\end{tabular}
\end{table}

\subsubsection*{৬. গভীর তাফসীর (\lat{Deep} \lat{Tafsir})}

\begin{itemize}
\item \textbf{ইবনে কাসীর:} ইবনে আব্বাস (রাঃ) বলেন — \lat{“}সাহূন\lat{”} মানে মুনাফিকরা, যারা নামাজ পড়ে জামাতে গেলে না, একা পড়ে যায়, আবার কখনো পড়ে না। সাঈদ ইবনে মুসাইয়িব বলেন — যে ব্যক্তি প্রথম ওয়াক্তে নামাজ পড়ে না, বরং শেষ সময়ে ছুটে ছুটে পড়ে — সে-ই উদাসীন। (তাফসীর ইবনে কাসীর, সূরা মাউন)
\item \textbf{কুরতুবী:} আয়াতের সতর্কবার্তা সেই নামাজিদের জন্য, যারা নামাজের সময়, রুকন বা শর্তের প্রতি অবহেলা করে — নামাজ পড়েও নামাজের প্রকৃত আদায় করে না। (তাফসীরুল কুরতুবী, সূরা মাউন)
\item \textbf{সা'দী:} \lat{“}উদাসীন\lat{”} অর্থ — নামাজের সময় ভুলে যাওয়া, নামাজ বিলম্বে পড়া, খুশুহীন নামাজ, রুকন ত্রুটিপূর্ণ রাখা। এ ধরনের নামাজই ব্যক্তিকে পাপ থেকে বিরত রাখে না। (আস-সা'দী, সূরা মাউন)
\end{itemize}
\subsubsection*{৭. সংশ্লিষ্ট হাদিস}

\begin{itemize}
\item \textbf{হাদিস (বুখারি ৫৫২, মুসলিম ৬২৬):} \lat{“}যে আসর ছেড়ে দিল, সে যেন পরিবার-ধন হারাল।\lat{”} — নামাজের অবহেলার ভয়াবহতা।
\item \textbf{হাদিস (বুখারি ৬৫৭, মুসলিম ৬৫১):} \lat{“}মুনাফিকদের উপর সবচেয়ে ভারী নামাজ — ইশা ও ফজর।\lat{”} — উদাসীনতার সাথে মুনাফিকির সম্পর্ক।
\end{itemize}
\subsubsection*{৮. গভীর শিক্ষা (\lat{Deep} \lat{Lessons})}

\begin{enumerate}
\item \textbf{নামাজ পড়েও অবহেলাকারী।} আয়াত সেই নামাজিদের সতর্ক করছে যারা পড়ে কিন্তু \textbf{\lat{negligent} (অবহেলা)} — অর্থাৎ নামাজ পড়া যথেষ্ট নয়, মানও চাই।
\item \textbf{সময়ের অবহেলা।} শেষ সময়ে ছুটে গিয়ে নামাজ পড়া — সাহূনের একটি রূপ।
\item \textbf{আত্মপরীক্ষা।} প্রতিটি মুমিনের প্রশ্ন — আমি কি নামাজে উদাসীন?
\end{enumerate}
\subsubsection*{৯. জীবনে প্রয়োগ (\lat{Practical} \lat{Application})}

\begin{enumerate}
\item নামাজের সময় হওয়ার সাথে সাথে আদায় করুন — দেরি করার \textbf{\lat{habit} (অভ্যাস)} ভাঙুন।
\item নামাজে খুশু ও রুকন-শর্ত ঠিক রাখার যত্ন নিন — শুধু \lat{“}পড়ে ফেলা\lat{”} নয়।
\end{enumerate}
\medskip\hrule\medskip

\subsection*{এক নজরে আরও আয়াত (সূত্র + অর্থ সংক্ষেপে)}

\subsubsection*{সূরা আল-জুমু'আ (৬২) — আয়াত ৯-১০}

\begin{center}\arab{يَا أَيُّهَا الَّذِينَ آمَنُوا إِذَا نُودِيَ لِلصَّلَاةِ مِن يَوْمِ الْجُمُعَةِ فَاسْعَوْا إِلَىٰ ذِكْرِ اللَّهِ وَذَرُوا الْبَيْعَ ۚ ذَٰلِكُمْ خَيْرٌ لَّكُمْ إِن كُنتُمْ تَعْلَمُونَ}\end{center}

\textbf{উচ্চারণ:} ইয়া-আইয়ুহাল্লাযীনা আ-মানূ ইযা- নূদিয়া লিস-সালা-তি মিং ইয়াওমিল জুমু'আতি ফাস'আও ইলা- যিকরিল্লা-হি ওয়া যারুল বাই'আ, যা-লিকুম খাইরুল্লাকুম ইন কুনতুম তা'লামূন।

\textbf{অর্থ (\lat{pure}):} \lat{“}হে ঈমানদারগণ! জুমার দিন নামাজের জন্য আহ্বান করা হলে আল্লাহর স্মরণের দিকে ধাবিত হও এবং ক্রয়-বিক্রয় ছেড়ে দাও। এটাই তোমাদের জন্য কল্যাণকর, যদি তোমরা জানতে।\lat{”}

\textbf{শিক্ষা:} জুমার আযানের পর ব্যবসা-বাণিজ্য বন্ধ করে মসজিদে যাওয়া ওয়াজিব — দুনিয়ার কাজের চেয়ে আল্লাহর স্মরণ অগ্রাধিকার পাবে।

\subsubsection*{সূরা আন-নিসা (৪) — আয়াত ১০১}

\begin{center}\arab{وَإِذَا ضَرَبْتُمْ فِي الْأَرْضِ فَلَيْسَ عَلَيْكُمْ جُنَاحٌ أَن تَقْصُرُوا مِنَ الصَّلَاةِ إِنْ خِفْتُمْ أَن يَفْتِنَكُمُ الَّذِينَ كَفَرُوا ۚ إِنَّ الْكَافِرِينَ كَانُوا لَكُمْ عَدُوًّا مُّبِينًا}\end{center}

\textbf{উচ্চারণ:} ওয়া ইযা- দ্বরাবতুম ফিল আরদ্বি ফালাইসা 'আলাইকুম জুনা-হুন আন তাক্বসুরূ মিনাস-সালা-তি ইন খিফতুম আই ইয়াফতিনাকুমুল্লাযীনা কাফারূ, ইন্নাল কা-ফিরীনা কা-নূ লাকুম 'আদুউয়াম মুবীনা।

\textbf{অর্থ (\lat{pure}):} \lat{“}আর যখন তোমরা পৃথিবীতে সফর করো, তখন নামাজ কসর (সংক্ষিপ্ত) করা তোমাদের জন্য অপরাধ নয় — যদি ভয় করো যে, কাফিররা তোমাদের ফিতনায় ফেলবে। নিশ্চয়ই কাফিররা তোমাদের প্রকাশ্য শত্রু।\lat{”}

\textbf{শিক্ষা:} সফরে চার রাকাতের নামাজ দুই রাকাত পড়ার বৈধতা কুরআনের এই আয়াত থেকে — কসরের ভিত্তি; হাদিসে নবীজির আমলও এটাই প্রতিষ্ঠা করেছে (বুখারি ১০৮২, মুসলিম ৬৮৬)।

\subsubsection*{সূরা আল-বাকারাহ (২) — আয়াত ১৫৩}

\begin{center}\arab{يَا أَيُّهَا الَّذِينَ آمَنُوا اسْتَعِينُوا بِالصَّبْرِ وَالصَّلَاةِ ۚ إِنَّ اللَّهَ مَعَ الصَّابِرِينَ}\end{center}

\textbf{উচ্চারণ:} ইয়া-আইয়ুহাল্লাযীনা আ-মানুসতা'ঈনূ বিস-সাবরি ওয়াস-সালা-হ, ইন্নাল্লা-হা মাআস-সা-বিরীন।

\textbf{অর্থ (\lat{pure}):} \lat{“}হে ঈমানদারগণ! ধৈর্য ও নামাজের মাধ্যমে সাহায্য চাও। নিশ্চয়ই আল্লাহ ধৈর্যশীলদের সাথে আছেন।\lat{”}

\textbf{শিক্ষা:} বাকারাহ ২:৪৫-এর পুনরাবৃত্তি মুমিনদের জন্য — বিপদে নামাজ ও ধৈর্যই আল্লাহর সাহায্যের পথ।

\medskip\hrule\medskip

\subsection*{সংক্ষিপ্ত সারাংশ (\lat{Summary} \lat{Table})}

\begin{table}[h]\centering\small
\begin{tabular}{ll}
আয়াত & মূল বার্তা \\
\hline
ত্বহা ২০:১৪ & নামাজের লক্ষ্য — আল্লাহর স্মরণ \\
বাকারাহ ২:৪৩ & নামাজ + জাকাত; জামাতের নির্দেশ \\
নিসা ৪:১০৩ & নামাজ নির্দিষ্ট সময়ে ফরজ \\
বাকারাহ ২:৪৫ & বিপদে ধৈর্য ও নামাজের সাহায্য \\
আনকাবুত ২৯:৪৫ & নামাজ ফাহেশা-মুনকার থেকে বিরত রাখে \\
মু'মিনুন ২৩:১-২ & সফলতার প্রথম গুণ — নামাজে খুশু \\
বাকারাহ ২:২৩৮ & সব নামাজে যত্ন, বিশেষত আসর \\
হুদ ১১:১১৪ & নেকি মন্দকে দূর করে; দিনের দুই প্রান্তে নামাজ \\
ইসরা ১৭:৭৮-৭৯ & সময়সীমা; ফজরের ফেরেশতা; তাহাজ্জুদ \\
মাউন ১০৭:৪-৫ & নামাজে উদাসীনতার ধিক্কার \\
জুমু'আ ৬২:৯ & জুমার আযানে কাজ ছেড়ে আল্লাহর দিকে \\
নিসা ৪:১০১ & সফরে কসরের বৈধতা \\
\hline
\end{tabular}
\end{table}


\end{document}
